\documentclass[11pt]{article}
\usepackage{mathblatt}
\def\mitzone{1}
\begin{document}
\blattkopf{Hypothesentests}{Gesamt}
\verzeichniszeile{\verz{zone}{Kennst du schon (Nr.~1–7)} \verztrenn \verz{e1}{1 Entscheidungsregel bestimmen: die Nullhypothese liefert p und die Binomialverteilung der Testgröße, die Alternative die Seite des Ablehnungsbereichs (rechtsseitig bei „mehr als“, linksseitig bei „weniger als“); die Grenze über kumulierte Wahrscheinlichkeiten so wählen, dass das Signifikanzniveau eingehalten wird (Nr.~8–10)} \verztrenn \verz{e2}{2 Die Nullhypothese wählen: aus der Sicht des Entscheiders (Nr.~11–12)} \verztrenn \verz{e3}{3 Fehlerarten und Güte: den Fehler zweiter Art für selbst gewählte Anteile berechnen (p dort wählen, wo die Nullhypothese falsch ist), einordnen und im Sachzusammenhang beschreiben; Mindestanteile für eine Fehlerschranke; die Gütekurve lesen (Ablehnwahrscheinlichkeit in Abhängigkeit von p (Nr.~13–15)} \verztrenn \verz{abhaken}{Das kann ich}}
% Zone „Das kennst du schon“ – aus bank/hypothesentests/zone.jsonl (zusammenbau v0.1)
\einheitenkopf[zone][Kennst du schon]{Das kennst du schon}
\setcounter{aufgabe}{0}

%% TODO Titel: Ich-kann-Satz fehlt in der Bank (Platzhalter: Kettenname) – Nr. 1
%% TODO Anweisung: ein Satz über der ersten Teilaufgabe fehlt in der Bank – Nr. 1
\begin{aufgabe}{Kumulierte Binomialwahrscheinlichkeiten berechnen und lesen}
\begin{teile}
\teil $X$ ist binomialverteilt mit $n = 3$ und $p = 0{,}5$. $P(X \le 1)$?
\teil $X$ ist binomialverteilt mit $n = 20$ und $p = 0{,}3$. $P(X \le 4)$ auf vier Stellen?
\end{teile}
\end{aufgabe}

%% TODO Titel: Ich-kann-Satz fehlt in der Bank (Platzhalter: Kettenname) – Nr. 2
%% TODO Anweisung: ein Satz über der ersten Teilaufgabe fehlt in der Bank – Nr. 2
\begin{aufgabe}{Gegenereignis und Komplementärsummen („mindestens“ gegen „höchstens“)}
\begin{teile}
\teil $P(X \le 5) = 0{,}8$. $P(X \ge 6)$?
\teil $X$ ist binomialverteilt mit $n = 50$ und $p = 0{,}2$; $P(X \le 12) \approx 0{,}8139$. $P(X \ge 13)$?
\end{teile}
\end{aufgabe}

%% TODO Titel: Ich-kann-Satz fehlt in der Bank (Platzhalter: Kettenname) – Nr. 3
%% TODO Anweisung: ein Satz über der ersten Teilaufgabe fehlt in der Bank – Nr. 3
\begin{aufgabe}{Erwartungswert n · p als Orientierung, auf welcher Seite der Ablehnungsbereich liegt}
\begin{teile}
\teil $n = 60$, $p = 0{,}5$: Erwartungswert $n \cdot p$?
\teil $n = 900$, $p = 0{,}035$: Erwartungswert $n \cdot p$?
\end{teile}
\end{aufgabe}

%% TODO Titel: Ich-kann-Satz fehlt in der Bank (Platzhalter: Kettenname) – Nr. 4
%% TODO Anweisung: ein Satz über der ersten Teilaufgabe fehlt in der Bank – Nr. 4
\begin{aufgabe}{Funktionsgraphen ablesen und deuten (Werte, Monotonie)}
\begin{teile}
\teil Lies $f(2)$ ab.

\begin{ksys}[xmin=0,xmax=8,ymin=0,ymax=5,ablesen]\funktion{0.5*\x+1}{f}\end{ksys}
\teil Lies ab: Für welches $x$ ist $f(x) = 1$?

\begin{ksys}[xmin=0,xmax=8,ymin=0,ymax=5,ablesen]\funktion{4/(\x+1)}{f}\end{ksys}
\end{teile}
\end{aufgabe}

%% TODO Titel: Ich-kann-Satz fehlt in der Bank (Platzhalter: Kettenname) – Nr. 5
%% TODO Anweisung: ein Satz über der ersten Teilaufgabe fehlt in der Bank – Nr. 5
\begin{aufgabe}{Den Rechner für kumulierte Verteilungswahrscheinlichkeiten einsetzen}
\begin{teile}
\teil $X$ ist binomialverteilt mit $n = 10$ und $p = 0{,}5$. $P(X \le 3)$ mit dem Rechner?
\teil $X$ ist binomialverteilt mit $n = 120$ und $p = 0{,}35$. $P(35 \le X \le 45)$?
\end{teile}
\end{aufgabe}

%% TODO Titel: Ich-kann-Satz fehlt in der Bank (Platzhalter: Kettenname) – Nr. 6
%% TODO Anweisung: ein Satz über der ersten Teilaufgabe fehlt in der Bank – Nr. 6
\begin{aufgabe}{Gegenereignis und Komplementärsummen („mindestens“ gegen „höchstens“) – Fehler finden}
\begin{teile}
\teil $X$ ist binomialverteilt mit $n = 30$ und $p = 0{,}2$. Tom rechnet: $P(X \ge 10) = 1 - P(X \le 10) \approx 1 - 0{,}9744 = 0{,}0256$. Finde den Fehler und rechne richtig.
\end{teile}
\end{aufgabe}

%% TODO Titel: Ich-kann-Satz fehlt in der Bank (Platzhalter: Kettenname) – Nr. 7
%% TODO Anweisung: ein Satz über der ersten Teilaufgabe fehlt in der Bank – Nr. 7
\begin{aufgabe}{Gegenereignis und Komplementärsummen („mindestens“ gegen „höchstens“) – selbst rechnen}
\begin{teile}
\teil $X$ ist binomialverteilt mit $n = 40$ und $p = 0{,}25$; $P(X \le 11) \approx 0{,}7151$. $P(X \ge 12)$?
\end{teile}
\end{aufgabe}

\clearpage
% Einheit 1 von 3 – Entscheidungsregel bestimmen: die Nullhypothese liefert p und die Binomialverteilung der Testgröße, die Alternative die Seite des Ablehnungsbereichs (rechtsseitig bei „mehr als“, linksseitig bei „weniger als“); die Grenze über kumulierte Wahrscheinlichkeiten so wählen, dass das Signifikanzniveau eingehalten wird (bank/hypothesentests/e1.jsonl, zusammenbau v0.1)
%% TODO Zeitmarke Einheit 1: Marken-Zeile nicht lesbar
%% TODO Zweigzeile Teil 1: was hier gelernt wird (Satz in Schülersprache aus dem Einheitstitel) – Einheit 1
\einheitenkopf[e1]{Einheit 1 von 3 · Entscheidungsregel bestimmen: die Nullhypothese liefert p und die Binomialverteilung der Testgröße, die Alternative die Seite des Ablehnungsbereichs (rechtsseitig bei „mehr als“, linksseitig bei „weniger als“); die Grenze über kumulierte Wahrscheinlichkeiten so wählen, dass das Signifikanzniveau eingehalten wird}
\zweigzeile{Abitur LK}
\setcounter{aufgabe}{7}

%% TODO Titel: Ich-kann-Satz fehlt in der Bank (Platzhalter: Kettenname) – Nr. 8
%% TODO Anweisung: ein Satz über der ersten Teilaufgabe fehlt in der Bank – Nr. 8
\begin{aufgabe}{Entscheidungsregel}
%% TODO \rechenplatz{n}: Schrittzahl des Grundfalls fehlt in der Bank (nur bei fester Schreibform, 2.3 b)
\begin{teile}
\teil $H_0\colon p \le 0{,}2$, Gegenhypothese $H_1\colon p > 0{,}2$. Auf welcher Seite liegt der Ablehnungsbereich? \\ \kreuz{rechtsseitig} \\ \kreuz{linksseitig}
\teil Eine Buchhändlerin vermutet, dass mehr als 25\,\% der Jugendlichen der Stadt Kunden in ihrem Laden sind; sonst plant sie eine Werbeaktion. Um unnötige Kosten zu vermeiden, soll die Nullhypothese $H_0$: „Höchstens 25\,\% der Jugendlichen sind Kunden“ mit 120 zufällig befragten Jugendlichen auf dem Signifikanzniveau 5\,\% getestet werden; $X$ ist die Anzahl der Kunden unter ihnen. Bestimme die Entscheidungsregel. \hfill (Abitur 2018 LK)
\teil Ein Reiseveranstalter verlängert ein Gewinnspiel nur, wenn mindestens 4\,\% der Teilnehmer danach eine Reise buchen. Getestet wird $H_0\colon p \ge 0{,}04$ ($p$ Buchungswahrscheinlichkeit) auf dem Signifikanzniveau 5\,\% mit 650 Teilnehmern; $X$ ist die Anzahl der Buchungen. Bestimme die Entscheidungsregel. \hfill (Abitur 2023 LK)
\teil Ein Fährbetrieb setzt eine zusätzliche Fähre für Radfahrer nur fort, wenn $H_0$: „Höchstens 12\,\% der Fahrgäste haben ein Rad dabei“ abgelehnt wird. Getestet wird auf dem Signifikanzniveau 8\,\% an 400 Fahrgästen; $X$ ist die Anzahl der Fahrgäste mit Rad. Bestimme die Entscheidungsregel und sichere die Grenze am Nachbarwert ab. \hfill (Abitur 2025 LK)
\teil Ein Streamingdienst testet $H_0$: „Höchstens 40\,\% der Nutzer sehen eine Serie zu Ende“ an 150 Nutzern auf dem Signifikanzniveau 5\,\%. Als Ablehnungsbereich wird $\{71; \ldots; 150\}$ angegeben, mit den Lösungsschritten: „$Y$ Anzahl der Nutzer, die die Serie zu Ende sehen; $P(Y \ge 71) \approx 0{,}041$.“ Begründe, warum diese Schritte nicht ausreichen, und ergänze sie. \hfill (Abitur 2024 LK)
\end{teile}
\end{aufgabe}

%% TODO Titel: Ich-kann-Satz fehlt in der Bank (Platzhalter: Kettenname) – Nr. 9
%% TODO Anweisung: ein Satz über der ersten Teilaufgabe fehlt in der Bank – Nr. 9
\begin{aufgabe}{Ablehnungsgrenze bei größerem Stichprobenumfang ohne höheren Fehler erster Art bestimmen}
\begin{teile}
\teil Ein Hersteller will höchstens 15\,\% fehlerhafte Teile. Bisher prüft er 80 Teile und bessert das Verfahren nach, wenn mindestens 18 fehlerhaft sind. Künftig prüft er 160 Teile; die Wahrscheinlichkeit für eine unnötige Nachbesserung soll dadurch nicht steigen. Ab wie vielen fehlerhaften Teilen wird nun nachgebessert? \hfill (Abitur 2017 LK)
\end{teile}
\end{aufgabe}

%% TODO Titel: Ich-kann-Satz fehlt in der Bank (Platzhalter: Kettenname) – Nr. 10
%% TODO Anweisung: ein Satz über der ersten Teilaufgabe fehlt in der Bank – Nr. 10
\begin{aufgabe}{Entscheidungsregel – Fehler finden, Begründen, Anwendung}
\begin{teile}
\teil Getestet wird $H_0$: „Mindestens 65\,\% der Kunden sind zufrieden“ an 90 Kunden auf dem Signifikanzniveau 5\,\%. Lisa bestimmt das kleinste $k$ mit $P(X \ge k) \le 0{,}05$ und erhält $k = 67$. Finde den Fehler.
\teil Begründe, warum bei $H_0\colon p \le 0{,}35$ gegen $H_1\colon p > 0{,}35$ der Ablehnungsbereich rechts liegt.
\teil Ein Kino kauft eine zweite Popcornmaschine nur, wenn $H_0$: „Höchstens 40\,\% der Besucher kaufen Popcorn“ abgelehnt wird; die Regel für 60 befragte Besucher lautet: Ablehnung ab 31 Käufern. In der Stichprobe kaufen 29. Was folgt daraus?
\end{teile}
\end{aufgabe}

\clearpage
% Einheit 2 von 3 – Die Nullhypothese wählen: aus der Sicht des Entscheiders (bank/hypothesentests/e2.jsonl, zusammenbau v0.1)
%% TODO Zeitmarke Einheit 2: Marken-Zeile nicht lesbar
%% TODO Zweigzeile Teil 1: was hier gelernt wird (Satz in Schülersprache aus dem Einheitstitel) – Einheit 2
\einheitenkopf[e2]{Einheit 2 von 3 · Die Nullhypothese wählen: aus der Sicht des Entscheiders}
\zweigzeile{Abitur LK}
\setcounter{aufgabe}{10}

%% TODO Titel: Ich-kann-Satz fehlt in der Bank (Platzhalter: Kettenname) – Nr. 11
%% TODO Anweisung: ein Satz über der ersten Teilaufgabe fehlt in der Bank – Nr. 11
\begin{aufgabe}{Wahl der Nullhypothese}
%% TODO \rechenplatz{n}: Schrittzahl des Grundfalls fehlt in der Bank (nur bei fester Schreibform, 2.3 b)
\begin{teile}
\teil $H_0$: „Das neue Medikament wirkt höchstens so oft wie das alte.“ In Wahrheit wirkt es nicht öfter, der Test lehnt $H_0$ aber ab. Welcher Fehler liegt vor? \\ \kreuz{Fehler erster Art} \\ \kreuz{Fehler zweiter Art} \\ \kreuz{kein Fehler}
\teil Bisher sind 5\,\% der Fahrradschläuche eines Herstellers undicht. Ein neues, teureres Gummi soll den Anteil senken. Getestet wird $H_0$: „Der Anteil undichter Schläuche beträgt mindestens 5\,\%.“ Gib eine Überlegung an, die zur Wahl dieser Nullhypothese geführt haben könnte, und begründe sie. \hfill (Abitur 2018 LK)
\teil Ein Verlag setzt bereits einen Algorithmus ein, der Artikel empfiehlt; bisher sind 65\,\% der Abonnenten zufrieden. Er soll nur abgeschaltet werden, wenn die Zufriedenheit nachweislich gesunken ist. Getestet wird $H_0$: „Mindestens 65\,\% der Abonnenten sind zufrieden.“ Gib eine mögliche Überlegung des Verlags zur Wahl dieser Nullhypothese an. \hfill (Abitur 2024 LK)
\teil Ein Händler hat mit einem Hersteller vereinbart, dass höchstens 5\,\% der gelieferten Ladegeräte defekt sein dürfen; sonst geht die Lieferung zurück. Er testet $H_0$: „Der Anteil defekter Geräte beträgt mindestens 5\,\%“ und nicht $H_0$: „Der Anteil beträgt höchstens 5\,\%.“ Gib eine Überlegung an, die zu dieser Wahl geführt haben könnte, und begründe sie. \hfill (Abitur 2022 LK)
\end{teile}
\end{aufgabe}

%% TODO Titel: Ich-kann-Satz fehlt in der Bank (Platzhalter: Kettenname) – Nr. 12
%% TODO Anweisung: ein Satz über der ersten Teilaufgabe fehlt in der Bank – Nr. 12
\begin{aufgabe}{Wahl der Nullhypothese – Fehler finden, Begründen, Anwendung}
\begin{teile}
\teil Mia begründet die Wahl von $H_0$: „Der Anteil defekter Geräte beträgt mindestens 5\,\%“ so: „Dann ist die Wahrscheinlichkeit, eine schlechte Lieferung anzunehmen, als Fehler zweiter Art begrenzt.“ Finde den Fehler.
\teil Begründe, warum nur der Fehler erster Art durch das Signifikanzniveau begrenzt ist.
\teil Eine Gemeinde will eine teure Lärmschutzwand nur bauen, wenn nachgewiesen ist, dass mehr als 25\,\% der Anwohner sich durch Verkehrslärm gestört fühlen. Wähle die Nullhypothese und begründe.
\end{teile}
\end{aufgabe}

\clearpage
% Einheit 3 von 3 – Fehlerarten und Güte: den Fehler zweiter Art für selbst gewählte Anteile berechnen (p dort wählen, wo die Nullhypothese falsch ist), einordnen und im Sachzusammenhang beschreiben; Mindestanteile für eine Fehlerschranke; die Gütekurve lesen (Ablehnwahrscheinlichkeit in Abhängigkeit von p (bank/hypothesentests/e3.jsonl, zusammenbau v0.1)
%% TODO Zeitmarke Einheit 3: Marken-Zeile nicht lesbar
%% TODO Zweigzeile Teil 1: was hier gelernt wird (Satz in Schülersprache aus dem Einheitstitel) – Einheit 3
\einheitenkopf[e3]{Einheit 3 von 3 · Fehlerarten und Güte: den Fehler zweiter Art für selbst gewählte Anteile berechnen (p dort wählen, wo die Nullhypothese falsch ist), einordnen und im Sachzusammenhang beschreiben; Mindestanteile für eine Fehlerschranke; die Gütekurve lesen (Ablehnwahrscheinlichkeit in Abhängigkeit von p}
\zweigzeile{Abitur LK}
\setcounter{aufgabe}{12}

%% TODO Titel: Ich-kann-Satz fehlt in der Bank (Platzhalter: Kettenname) – Nr. 13
%% TODO Anweisung: ein Satz über der ersten Teilaufgabe fehlt in der Bank – Nr. 13
\begin{aufgabe}{Fehlerarten und Güte}
%% TODO \rechenplatz{n}: Schrittzahl des Grundfalls fehlt in der Bank (nur bei fester Schreibform, 2.3 b)
\begin{teile}
\teil Getestet wird $H_0\colon p \le 0{,}25$. Für welchen wahren Anteil $p$ ist ein Fehler zweiter Art möglich? \\ \kreuz{$p = 0{,}2$} \\ \kreuz{$p = 0{,}25$} \\ \kreuz{$p = 0{,}35$}
\teil Ein Reiseportal verlängert ein Gewinnspiel, wenn sich das lohnt; das ist bei einer Buchungswahrscheinlichkeit $p$ von mindestens 4\,\% der Fall. Es testet $H_0\colon p \ge 0{,}04$ mit 500 Teilnehmern und lehnt $H_0$ ab, wenn weniger als 13 buchen. Berechne den Fehler zweiter Art für zwei geeignete Werte von $p$ und deute ihn im Sachzusammenhang. \hfill (Abitur 2023 LK)
\teil Ein Weinkenner behauptet, er erkenne den Jahrgang am Geschmack. Bei 12 Proben wird $H_0$: „Er entscheidet höchstens mit Wahrscheinlichkeit 0,45 richtig“ ab 9 richtigen Antworten abgelehnt. Die Abbildung zeigt die Wahrscheinlichkeit, $H_0$ abzulehnen, in Abhängigkeit von der Trefferwahrscheinlichkeit $p$. Wähle ein geeignetes $p$ und bestimme mithilfe der Abbildung den Fehler zweiter Art. \hfill (Abitur 2018 LK)

\begin{ksys}[xmin=0,xmax=1,ymin=0,ymax=1,xstep=0.1,ystep=0.1,ablesen,xlabel=$p$]\funktion{220*\x^9*(1-\x)^3+66*\x^10*(1-\x)^2+12*\x^11*(1-\x)+\x^12}{}\end{ksys}
\teil Getestet wird $H_0$: „Höchstens 35\,\% der Besucher kommen mit dem Rad“ an 300 Besuchern; der Ablehnungsbereich ist $\{120; \ldots; 300\}$. Weise nach, dass der Fehler zweiter Art mehr als 90\,\% betragen kann. \hfill (Abitur 2024 LK)
\teil Ein Fährbetrieb testet $H_0$: „Höchstens 12\,\% der Fahrgäste haben ein Rad dabei“ an 300 Fahrgästen und lehnt ab, wenn mehr als 45 ein Rad dabeihaben; nur dann fährt eine zusätzliche Radfähre. Der Fehler zweiter Art soll höchstens 10\,\% betragen. Wie groß muss der Anteil der Fahrgäste mit Rad mindestens sein (auf ganze Prozent)? Beschreibe den Fehler zweiter Art im Sachzusammenhang. \hfill (Abitur 2025 LK)
\teil Ein Händler testet $H_0$: „Mindestens 5\,\% der Armbänder sind fehlerhaft“; abgelehnt wird bei höchstens 3 fehlerhaften Armbändern in der Stichprobe. Die Abbildung zeigt für den gewählten Stichprobenumfang die Wahrscheinlichkeit, $H_0$ abzulehnen, in Abhängigkeit vom Anteil $p$. Begründe mithilfe der Abbildung, dass der Stichprobenumfang größer als hundert ist. \hfill (Abitur 2022 LK)

\begin{ksys}[xmin=0,xmax=0.1,ymin=0,ymax=1,xstep=0.01,ystep=0.1,ablesen,xlabel=$p$]\funktionab{(1-\x)^150+exp(5.0106+1*ln(\x)+149*ln(1-\x))+exp(9.3214+2*ln(\x)+148*ln(1-\x))+exp(13.2200+3*ln(\x)+147*ln(1-\x))}{}{0.001}{0.1}\end{ksys}
\end{teile}
\end{aufgabe}

%% TODO Titel: Ich-kann-Satz fehlt in der Bank (Platzhalter: Kettenname) – Nr. 14
%% TODO Anweisung: ein Satz über der ersten Teilaufgabe fehlt in der Bank – Nr. 14
\begin{aufgabe}{Fehlentscheidungen eines Tests im Sachzusammenhang beschreiben}
\begin{teile}
\teil Ein Hersteller von LED-Lampen will, dass höchstens 8\,\% der Lampen früh ausfallen. Er prüft 120 Lampen und stellt die Fertigung nur um, wenn mindestens 15 früh ausfallen. Beschreibe die Fehlentscheidungen, die bei dieser Kontrolle auftreten können. \hfill (Abitur 2017 LK)
\end{teile}
\end{aufgabe}

%% TODO Titel: Ich-kann-Satz fehlt in der Bank (Platzhalter: Kettenname) – Nr. 15
%% TODO Anweisung: ein Satz über der ersten Teilaufgabe fehlt in der Bank – Nr. 15
\begin{aufgabe}{Fehlerarten und Güte – Fehler finden, Begründen, Anwendung, Darstellungswechsel}
\begin{teile}
\teil Getestet wird $H_0\colon p \le 0{,}35$ mit 40 Versuchen, abgelehnt wird ab 20 Treffern. Tom berechnet den Fehler zweiter Art für $p = 0{,}3$: $P(X \le 19) \approx 0{,}9937$. Finde den Fehler.
\teil Begründe, warum der Fehler zweiter Art vom gewählten wahren Anteil $p$ abhängt.
\teil Eine Abfüllanlage wird gewartet, wenn der Test von $H_0$: „Höchstens 2,5\,\% der Flaschen sind falsch befüllt“ bei 300 geprüften Flaschen ablehnt, also ab 13 falsch befüllten. Wie wahrscheinlich bleibt eine Anlage mit 5\,\% Fehlbefüllung ungewartet? Ist das vertretbar?
\teil Getestet wird $H_0\colon p \le 0{,}4$ mit 15 Versuchen, abgelehnt wird ab 10 Treffern. Lies die Ablehnwahrscheinlichkeiten für die $p$-Werte der Tabelle an der Gütekurve ab und trage sie ein.

\begin{ksys}[xmin=0,xmax=1,ymin=0,ymax=1,xstep=0.1,ystep=0.1,ablesen,xlabel=$p$]\funktion{3003*\x^10*(1-\x)^5+1365*\x^11*(1-\x)^4+455*\x^12*(1-\x)^3+105*\x^13*(1-\x)^2+15*\x^14*(1-\x)+\x^15}{}\end{ksys} \sachtabelle{lcccc}{$p$ & 0,6 & 0,7 & 0,8 & 0,9}{Ablehnwahrscheinlichkeit & \leerzelle & \leerzelle & \leerzelle & \leerzelle}
\end{teile}
\end{aufgabe}

% Abhakseite „Das kann ich“ – Zone nur im Gesamt (\mitzone)
%% TODO Ich-kann-Sätze: Platzhalter wie in den Aufgabendateien
\begin{abhakseite}
\ifdefined\mitzone
\abhakgruppe{Kennst du schon}
\abhak{1}{Kumulierte Binomialwahrscheinlichkeiten berechnen und lesen}
\abhak{2}{Gegenereignis und Komplementärsummen („mindestens“ gegen „höchstens“)}
\abhak{3}{Erwartungswert n · p als Orientierung, auf welcher Seite der Ablehnungsbereich liegt}
\abhak{4}{Funktionsgraphen ablesen und deuten (Werte, Monotonie)}
\abhak{5}{Den Rechner für kumulierte Verteilungswahrscheinlichkeiten einsetzen}
\abhak{6}{Gegenereignis und Komplementärsummen („mindestens“ gegen „höchstens“) – Fehler finden}
\abhak{7}{Gegenereignis und Komplementärsummen („mindestens“ gegen „höchstens“) – selbst rechnen}
\fi
\abhakgruppe{Einheit 1 · Entscheidungsregel bestimmen: die Nullhypothese liefert p und die Binomialverteilung der Testgröße, die Alternative die Seite des Ablehnungsbereichs (rechtsseitig bei „mehr als“, linksseitig bei „weniger als“); die Grenze über kumulierte Wahrscheinlichkeiten so wählen, dass das Signifikanzniveau eingehalten wird}
\abhak{8}{Entscheidungsregel}
\abhak{9}{Ablehnungsgrenze bei größerem Stichprobenumfang ohne höheren Fehler erster Art bestimmen}
\abhak{10}{Entscheidungsregel – Fehler finden, Begründen, Anwendung}
\abhakgruppe{Einheit 2 · Die Nullhypothese wählen: aus der Sicht des Entscheiders}
\abhak{11}{Wahl der Nullhypothese}
\abhak{12}{Wahl der Nullhypothese – Fehler finden, Begründen, Anwendung}
\abhakgruppe{Einheit 3 · Fehlerarten und Güte: den Fehler zweiter Art für selbst gewählte Anteile berechnen (p dort wählen, wo die Nullhypothese falsch ist), einordnen und im Sachzusammenhang beschreiben; Mindestanteile für eine Fehlerschranke; die Gütekurve lesen (Ablehnwahrscheinlichkeit in Abhängigkeit von p}
\abhak{13}{Fehlerarten und Güte}
\abhak{14}{Fehlentscheidungen eines Tests im Sachzusammenhang beschreiben}
\abhak{15}{Fehlerarten und Güte – Fehler finden, Begründen, Anwendung, Darstellungswechsel}
\end{abhakseite}

\end{document}
